% Closed table for Chapter 6.
% Usage:
%   % Closed table for Chapter 6.
% Usage:
%   % Closed table for Chapter 6.
% Usage:
%   % Closed table for Chapter 6.
% Usage:
%   \input{Figures/06/table_application_module_roles.tex}
%   \ApplicationModuleTable[0.95\textwidth]
%
% Change the optional width argument to resize the complete table.
% The table uses tabularx and array.

\newcommand{\ApplicationModuleTable}[1][0.97\textwidth]{%
\begin{table}[H]
    \centering
    \caption{Main software groups and their functional responsibilities.}
    \label{tab:application_module_roles}
    \renewcommand{\arraystretch}{1.18}
    \setlength{\tabcolsep}{6pt}
    \footnotesize
    \begin{tabularx}{#1}{
        |>{\raggedright\arraybackslash\hsize=0.60\hsize}X
        |>{\raggedright\arraybackslash\hsize=1.40\hsize}X|
    }
        \hline
        \textbf{Module or directory} &
        \textbf{Main responsibility} \\
        \hline

        \path{app_interactiva.py} &
        Coordinates file inputs, loaded cases and execution of the analysis tabs. \\
        \hline

        \path{core/thermaldata.py} &
        Reads each result file and provides access to nodes, attributes,
        time histories and conductor data. \\
        \hline

        \path{app_helpers.py}, \path{viewer_utils.py} &
        Provide hierarchy management, selection validation, unit handling,
        case-comparison operations and common output preparation. \\
        \hline

        \path{tabs/} &
        Contains one module for each analysis tab, including case inspection,
        heat-flow exploration, conductor analysis, transient post-processing,
        and the electrical balance and battery model. \\
        \hline

        \path{scene_manager.py}, \path{scene_ui.py},
        \path{restored_scenes.py} &
        Provide the controls required to save, list, restore and export complete
        analysis configurations and their generated outputs. \\
        \hline
    \end{tabularx}
\end{table}%
}

%   \ApplicationModuleTable[0.95\textwidth]
%
% Change the optional width argument to resize the complete table.
% The table uses tabularx and array.

\newcommand{\ApplicationModuleTable}[1][0.97\textwidth]{%
\begin{table}[H]
    \centering
    \caption{Main software groups and their functional responsibilities.}
    \label{tab:application_module_roles}
    \renewcommand{\arraystretch}{1.18}
    \setlength{\tabcolsep}{6pt}
    \footnotesize
    \begin{tabularx}{#1}{
        |>{\raggedright\arraybackslash\hsize=0.60\hsize}X
        |>{\raggedright\arraybackslash\hsize=1.40\hsize}X|
    }
        \hline
        \textbf{Module or directory} &
        \textbf{Main responsibility} \\
        \hline

        \path{app_interactiva.py} &
        Coordinates file inputs, loaded cases and execution of the analysis tabs. \\
        \hline

        \path{core/thermaldata.py} &
        Reads each result file and provides access to nodes, attributes,
        time histories and conductor data. \\
        \hline

        \path{app_helpers.py}, \path{viewer_utils.py} &
        Provide hierarchy management, selection validation, unit handling,
        case-comparison operations and common output preparation. \\
        \hline

        \path{tabs/} &
        Contains one module for each analysis tab, including case inspection,
        heat-flow exploration, conductor analysis, transient post-processing,
        and the electrical balance and battery model. \\
        \hline

        \path{scene_manager.py}, \path{scene_ui.py},
        \path{restored_scenes.py} &
        Provide the controls required to save, list, restore and export complete
        analysis configurations and their generated outputs. \\
        \hline
    \end{tabularx}
\end{table}%
}

%   \ApplicationModuleTable[0.95\textwidth]
%
% Change the optional width argument to resize the complete table.
% The table uses tabularx and array.

\newcommand{\ApplicationModuleTable}[1][0.97\textwidth]{%
\begin{table}[H]
    \centering
    \caption{Main software groups and their functional responsibilities.}
    \label{tab:application_module_roles}
    \renewcommand{\arraystretch}{1.18}
    \setlength{\tabcolsep}{6pt}
    \footnotesize
    \begin{tabularx}{#1}{
        |>{\raggedright\arraybackslash\hsize=0.60\hsize}X
        |>{\raggedright\arraybackslash\hsize=1.40\hsize}X|
    }
        \hline
        \textbf{Module or directory} &
        \textbf{Main responsibility} \\
        \hline

        \path{app_interactiva.py} &
        Coordinates file inputs, loaded cases and execution of the analysis tabs. \\
        \hline

        \path{core/thermaldata.py} &
        Reads each result file and provides access to nodes, attributes,
        time histories and conductor data. \\
        \hline

        \path{app_helpers.py}, \path{viewer_utils.py} &
        Provide hierarchy management, selection validation, unit handling,
        case-comparison operations and common output preparation. \\
        \hline

        \path{tabs/} &
        Contains one module for each analysis tab, including case inspection,
        heat-flow exploration, conductor analysis, transient post-processing,
        and the electrical balance and battery model. \\
        \hline

        \path{scene_manager.py}, \path{scene_ui.py},
        \path{restored_scenes.py} &
        Provide the controls required to save, list, restore and export complete
        analysis configurations and their generated outputs. \\
        \hline
    \end{tabularx}
\end{table}%
}

%   \ApplicationModuleTable[0.95\textwidth]
%
% Change the optional width argument to resize the complete table.
% The table uses tabularx and array.

\newcommand{\ApplicationModuleTable}[1][0.97\textwidth]{%
\begin{table}[H]
    \centering
    \caption{Main software groups and their functional responsibilities.}
    \label{tab:application_module_roles}
    \renewcommand{\arraystretch}{1.18}
    \setlength{\tabcolsep}{6pt}
    \footnotesize
    \begin{tabularx}{#1}{
        |>{\raggedright\arraybackslash\hsize=0.60\hsize}X
        |>{\raggedright\arraybackslash\hsize=1.40\hsize}X|
    }
        \hline
        \textbf{Module or directory} &
        \textbf{Main responsibility} \\
        \hline

        \path{app_interactiva.py} &
        Coordinates file inputs, loaded cases and execution of the analysis tabs. \\
        \hline

        \path{core/thermaldata.py} &
        Reads each result file and provides access to nodes, attributes,
        time histories and conductor data. \\
        \hline

        \path{app_helpers.py}, \path{viewer_utils.py} &
        Provide hierarchy management, selection validation, unit handling,
        case-comparison operations and common output preparation. \\
        \hline

        \path{tabs/} &
        Contains one module for each analysis tab, including case inspection,
        heat-flow exploration, conductor analysis, transient post-processing,
        and the electrical balance and battery model. \\
        \hline

        \path{scene_manager.py}, \path{scene_ui.py},
        \path{restored_scenes.py} &
        Provide the controls required to save, list, restore and export complete
        analysis configurations and their generated outputs. \\
        \hline
    \end{tabularx}
\end{table}%
}
